\chapter*{Conclusion et perspectives}
\addcontentsline{toc}{chapter}{Conclusion et perspectives}
\markboth{Conclusion et perspectives}{}

Ce projet avait pour objectif de concevoir, avec des composants à faible coût, un système IoT capable de détecter les
obstacles sur le trajet d'un véhicule, d'en évaluer la gravité et d'exploiter ces informations en temps réel ; au terme
de ce travail, une chaîne complète a été réalisée, du capteur jusqu'à la décision. Sur le plan matériel, nous avons
construit une voiture miniature autour d'une carte Arduino UNO, d'un pont en H L298N, de deux moteurs à courant
continu, de deux capteurs infrarouges et d'un capteur ultrasonique, qui suit de façon autonome une piste tracée au
ruban adhésif et s'arrête devant un obstacle, tandis qu'une caméra ESP32-CAM diffuse en Wi-Fi les images prises de son
point de vue. Sur le plan de l'intelligence artificielle, nous avons entraîné des détecteurs YOLO26 dans un
environnement Docker reproductible : le détecteur d'anomalies routières, affiné pendant huit heures sur 7\,894 images
réelles, atteint un mAP@50 de 0,53 en temps réel (0,57 au meilleur réglage) sur des lieux absents de l'entraînement,
et le modèle du prototype, entraîné sur notre propre jeu de données annoté automatiquement avec SAM~2, reconnaît tous
les obstacles de son ensemble de test et détecte la boîte sur le flux réel de la caméra avec une confiance de
95\,\%. Sur le plan logiciel, l'application \app, développée avec Dash et Plotly, exécute ces modèles en temps réel,
transforme leurs détections en obstacles localisés et priorisés selon un score explicable, les intègre à la
planification des itinéraires et permet de superviser avec la même chaîne logicielle une flotte sur le réseau routier
réel d'un quartier de Paris et la voiture du prototype. Ce projet nous a permis de mettre en pratique l'ensemble des
notions du module (capteurs et actionneurs, microcontrôleurs, communication sans fil, traitement des données,
intelligence artificielle et supervision) et d'apprendre à évaluer honnêtement un modèle en identifiant ses limites.
Plusieurs perspectives permettraient de prolonger ce travail : enrichir les données du prototype avec des images
enregistrées par l'ESP32-CAM elle-même, dans différents lieux et éclairages, pour supprimer les fausses détections sur
des objets inconnus ; améliorer le modèle routier par un entraînement plus long et davantage d'exemples de
nids-de-poule ; relier l'Arduino au système pour transmettre la distance ultrasonique et recevoir des consignes du
serveur ; localiser précisément la voiture grâce aux marqueurs ArUco sur la piste et, dans le monde réel, à un module
GPS ; exécuter les modèles sur un calculateur embarqué (Raspberry Pi, Jetson) pour réduire la latence ; et enfin
déployer l'application sur un serveur accessible à distance, avec des notifications en cas d'obstacle critique.
